\section{Request configurations and reproducibility}
\label{app:icml-reproducibility}

\subsection{What the model labels identify}
All 1,008 requests use the gateway endpoint
\nolinkurl{https://openrouter.ai/api/v1/images}.
Table~\ref{tab:icml-request-configurations} records the exact requested model
identifiers and provider restrictions. Every request sets
\texttt{n=1}, \texttt{aspect\_ratio=1:1}, and
\texttt{provider.allow\_fallbacks=false}; the listed provider is the sole
entry in \texttt{provider.only}. The payloads contain no image input or seed
field. Separate requests supply the two repeats.

\begin{table*}[t]
\caption{Recorded request configurations for the six-model experiment.
Rendering entries are explicit payload fields, not independently verified
service settings. All 1,008 returned images decode to $1024\times1024$ pixels.
None of the recorded responses echoes a model, quality, size or output-format
field.}
\label{tab:icml-request-configurations}
\centering\small
\setlength{\tabcolsep}{5pt}
\begin{tabularx}{\textwidth}{@{}>{\raggedright\arraybackslash}p{.39\textwidth}>{\raggedright\arraybackslash}p{.24\textwidth}>{\raggedright\arraybackslash}X@{}}
\toprule
Requested model identifier & Sole requested provider & Additional rendering fields\\
\midrule
\nolinkurl{openai/gpt-image-1} & \texttt{openai} & \texttt{quality=medium}; \texttt{background=opaque}\\[3pt]
\nolinkurl{openai/gpt-image-2} & \texttt{openai} & \texttt{quality=medium}; \texttt{background=opaque}\\[3pt]
\nolinkurl{openai/gpt-image-2.5-flare} & \texttt{openai} & \texttt{quality=medium}; \texttt{background=opaque}\\[3pt]
\nolinkurl{openai/gpt-image-2.5-sunburst} & \texttt{openai} & \texttt{quality=medium}; \texttt{background=opaque}\\[3pt]
\nolinkurl{google/gemini-3.1-flash-image} & \texttt{google-ai-studio} & \texttt{resolution=1K}\\[3pt]
\nolinkurl{black-forest-labs/flux.2-max} & \nolinkurl{black-forest-labs/us-3} & \texttt{output\_format=png}\\
\bottomrule
\end{tabularx}
\end{table*}

These records establish the configuration submitted to the gateway, not the
identity of hidden served weights. Provider catalogue records are bound to the
precollection freeze, but the absent echoed fields prevent response-level
confirmation of model and quality. The labels therefore denote requested
configurations. The earlier interface's acceptance of an invented model name
motivated this explicit-route collection; it does not prove that valid names
selected identical models. The terminated predecessor and technical probes
are excluded from the 1,008-image analysis.

\subsection{Timeline and evidence bindings}
The run identifier is \texttt{psv2-20260911}. Its local calendar date is
11 September 2026 (UTC+9), whereas the collection occurred on 10 September
in UTC. The freeze was recorded at 15:34:19.942590 UTC, the first request
started at 15:34:20.846233 UTC, and the terminal collection receipt was
recorded at 18:21:04.039850 UTC. The measurement receipt follows at
18:30:03.866908 UTC. Thus the date in the protocol and run identifier does
not imply a freeze after collection.

The freeze binds 27 files, including the protocols, exact assignments,
provider catalogue records, implementation, numerical tests, reference
features, development scaler and locked dependencies. The terminal record
contains 1,008 successes in 1,008 attempts, with a response hash and image
hash for every output. These hashes support consistency checks against the
retained bytes; local timestamps and hashes are not an independent timestamping
or preregistration service. A separate reader correction selects the intended
649 measured references from a manifest retaining four historical failed
records; it does not recollect those records.

\subsection{Numerical replay and image access}
The compact artifact supports numerical replay from retained feature vectors,
scalers, manifests and versioned results. This covers the four primary views
(full-frame and square-window measurements, each with the original or
content-reweighted reference target), both retrospective diagnostic versions,
and the source-region/scaler/label sensitivities. These checks verify the
recorded computation conditional on those inputs. They do not independently
validate the feature extraction, the source judgments, or served model identity.

Exact-pixel verification additionally requires retained local source files.
The terminal audit verifies the current experiment's response and image
bytes; the example check verifies the 40 displayed sources. The source-quality
image check verifies all 870 reference/development source hashes and
re-extracts features for the 131 audited crops. Exact source pixels remain
outside the compact repository, and the current experiment lacks a dedicated
public exact-pixel archive or verified recovery route. Numerical reproducibility
is therefore more complete than independent image-measurement reproducibility.

\subsection{Verification commands}
Run from the repository root using the locked Python 3.13.11 environment and
\texttt{uv}. The following targets are offline and do not generate images:
\begin{quote}\small
\begin{verbatim}
make specificity-check
make review-check
make reference-quality-check
make figures-check
\end{verbatim}
\end{quote}
The first three targets replay scientific results; the last checks the
manuscript's numerical figures and generated tables. Focused tests for the
primary estimator, routes, corrected reader and diagnostic arithmetic are:
\begin{quote}\small
\begin{verbatim}
uv run --locked pytest -q tests/painter_specificity_v1 \
  tests/painter_specificity_v2 \
  tests/painter_specificity_measurement_v1 \
  tests/painter_specificity_review_v1 \
  tests/painter_specificity_review_v2 \
  tests/painter_reference_quality_v1
\end{verbatim}
\end{quote}
With the retained source pixels and responses available, run:
\begin{quote}\small
\begin{verbatim}
make specificity-audit
make example-images-check
make reference-quality-images-check
\end{verbatim}
\end{quote}
These image-dependent commands add byte verification or feature re-extraction;
they do not collect new outputs. None of these checks is a scientific peer-review
rating or a validation of perceptual artistic fidelity.
